\documentclass[11pt,a4paper]{article}

\usepackage[utf8]{inputenc}
\usepackage[T1]{fontenc}
\usepackage[margin=2.5cm]{geometry}
\usepackage{hyperref}
\usepackage{listings}
\usepackage{xcolor}
\usepackage{titlesec}
\usepackage{enumitem}
\usepackage{booktabs}
\usepackage{parskip}

% Colors
\definecolor{codebg}{RGB}{245,245,245}
\definecolor{codeframe}{RGB}{200,200,200}
\definecolor{keyword}{RGB}{0,0,180}
\definecolor{comment}{RGB}{80,120,80}
\definecolor{string}{RGB}{160,0,0}

% Code listing style
\lstdefinestyle{code}{
  backgroundcolor=\color{codebg},
  frame=single,
  rulecolor=\color{codeframe},
  basicstyle=\ttfamily\small,
  keywordstyle=\color{keyword}\bfseries,
  commentstyle=\color{comment}\itshape,
  stringstyle=\color{string},
  breaklines=true,
  breakatwhitespace=true,
  tabsize=4,
  showstringspaces=false,
  numbers=none,
}

\lstdefinestyle{yaml}{
  style=code,
  language=,
  morekeywords={management,tracing,sampling,probability,otlp,endpoint},
}

\lstdefinestyle{kotlin}{
  style=code,
  language=Java,
  morekeywords={fun,val,var,override,data,object,companion,when,is,in,by,
                import,package,class,interface,object,open,private,internal,
                suspend,inline,reified,sealed,enum,annotation,true,false,null},
}

\lstdefinestyle{xml}{
  style=code,
  language=XML,
}

% Section formatting
\titleformat{\section}{\large\bfseries}{}{0em}{}[\titlerule]
\titleformat{\subsection}{\normalsize\bfseries}{}{0em}{}
\titleformat{\subsubsection}{\normalsize\itshape}{}{0em}{}

\hypersetup{
  colorlinks=true,
  linkcolor=blue,
  urlcolor=blue,
}

\title{\textbf{Observability Changes}\\\large Distributed Tracing --- final-project-g11}
\author{Session notes}
\date{April 2026}

\begin{document}

\maketitle
\tableofcontents
\newpage

% -----------------------------------------------------------------------------
\section{Overview}

This document records every change made to introduce distributed tracing across
all four microservices of the project, propagate trace context through the
Debezium outbox so the entire upload pipeline appears as a single trace, and fix
the AI worker file-download bug uncovered in the process.

\medskip
\begin{center}
\begin{tabular}{lll}
\toprule
\textbf{Service} & \textbf{Spring Boot} & \textbf{Port} \\
\midrule
TusUploaderService  & 3.5.8  & 8081 \\
AIProcessingService & 3.5.9  & 8083 \\
ApiGateway          & 3.4.12 & 8080 \\
UserDetailService   & 3.5.8  & 8082 \\
\bottomrule
\end{tabular}
\end{center}

\medskip
The observability backend is the \textbf{Grafana LGTM} all-in-one container
(\texttt{grafana/otel-lgtm}) already present in \texttt{docker-compose.yml},
exposing:
\begin{itemize}[noitemsep]
  \item \texttt{localhost:3000} --- Grafana UI
  \item \texttt{localhost:3100} --- Loki (log ingestion, used by the Loki4j appender)
  \item \texttt{localhost:4318} --- OTLP HTTP (traces $\to$ Tempo)
  \item \texttt{localhost:9099} --- Prometheus (metrics)
\end{itemize}

The expected end-to-end trace tree for a single file upload is:

\begin{lstlisting}[style=code, caption={Trace tree visible in Tempo after all changes}]
HTTP POST /upload  (TusUploaderService)
  └── kafka receive file-outbox  (AIProcessingService: MetadataProcessor)
        └── updateFileMetadata  (TusUploaderService: KafkaConsumer)
        └── kafka receive file-analyzed  (AIProcessingService: EmbeddingProcessor)
              └── embeddings-batch publish x N
\end{lstlisting}

% -----------------------------------------------------------------------------
\section{Tracing Dependencies --- All 4 Services}

\subsection{Why Micrometer Tracing instead of the OTel Java Agent}

All four \texttt{build.gradle.kts} files previously contained a
\texttt{downloadOtelAgent} Gradle task that fetched
\texttt{opentelemetry-javaagent.jar} v2.12.0 and attached it to
\texttt{bootRun} via a \texttt{-javaagent} JVM argument.

This approach was \textbf{removed} because:
\begin{itemize}
  \item Running both the OTel Java agent and Spring Boot's
        \texttt{micrometer-tracing-bridge-otel} at the same time produces
        \textbf{duplicate spans} for every operation.
  \item The logback MDC keys already in use (\texttt{traceId}, \texttt{spanId})
        match the \textbf{Micrometer Tracing} naming convention, not the OTel
        agent convention (\texttt{trace\_id}, \texttt{span\_id}).
  \item The Spring Boot native approach integrates cleanly with Kafka
        instrumentation and the rest of the Spring ecosystem.
\end{itemize}

\subsection{Dependencies added}

In each \texttt{build.gradle.kts}, under the \texttt{dependencies} block:

\begin{lstlisting}[style=kotlin, caption={build.gradle.kts --- added to all 4 services}]
// Distributed Tracing (Micrometer + OTel bridge + OTLP exporter)
implementation("io.micrometer:micrometer-tracing-bridge-otel")
implementation("io.opentelemetry:opentelemetry-exporter-otlp")
\end{lstlisting}

Versions are managed by the Spring Boot BOM.
\texttt{micrometer-tracing-bridge-otel} bridges Spring's \texttt{Tracer}
abstraction to the OTel SDK; \texttt{opentelemetry-exporter-otlp} sends spans
to Tempo via OTLP HTTP.  The bridge also auto-configures a
\texttt{Propagator} bean (W3C TraceContext) that is used for manual context
extraction described in Section~\ref{sec:debezium}.

\subsection{Configuration added to application.yml}

\begin{lstlisting}[style=yaml, caption={application.yml --- added to all 4 services}]
management:
  tracing:
    sampling:
      probability: 1.0        # sample every request (tune down for production)
  otlp:
    tracing:
      endpoint: http://localhost:4318/v1/traces
\end{lstlisting}

\subsection{Console appender --- logback-spring.xml}

The Loki appender pattern already contained \texttt{traceId=\%mdc\{traceId:-\}}
in all services.  The \textbf{CONSOLE} appender was missing it. Updated in all
four \texttt{logback-spring.xml} files:

\begin{lstlisting}[style=xml, caption={logback-spring.xml --- CONSOLE appender (all 4 services)}]
<appender name="CONSOLE" class="ch.qos.logback.core.ConsoleAppender">
    <encoder>
        <pattern>
            ${LOG_DATE} ${LOG_THREAD} ${LOG_LEVEL} ${LOG_LOGGER}
            %clr([%mdc{traceId:-} %mdc{spanId:-}]){faint}
            --> %msg%n%ex{full}
        </pattern>
    </encoder>
</appender>
\end{lstlisting}

% -----------------------------------------------------------------------------
\section{AIProcessingService --- Kafka Consumer Tracing}

\subsection{Problem}

\texttt{KafkaConfig.kt} manually constructs both \texttt{ConsumerFactory} and
\texttt{ConcurrentKafkaListenerContainerFactory} from scratch, completely
bypassing Spring Boot's auto-configuration.  Without
\texttt{isObservationEnabled = true}, no Kafka consumer spans were ever created,
making \texttt{AIProcessingService} invisible in Tempo.

\subsection{Two-factory design}

AIProcessingService has two distinct consumers:

\begin{itemize}
  \item \textbf{MetadataProcessor} listens to \texttt{file-outbox} --- a
        Debezium CDC topic.  Debezium is the actual Kafka publisher and never
        sets W3C trace headers, so Spring Kafka's built-in observation would
        create \textbf{orphan root spans} for every message.  Trace context is
        extracted manually inside the listener (see Section~\ref{sec:debezium}).

  \item \textbf{EmbeddingProcessor} listens to \texttt{file-analyzed} --- a
        topic published by \texttt{MetadataProcessor} via \texttt{KafkaTemplate}
        with explicit \texttt{traceparent} headers.  Spring Kafka's observation
        correctly creates a child span here.
\end{itemize}

These two cases require different configuration, so two factory beans are
defined in \texttt{KafkaConfig.kt}:

\begin{lstlisting}[style=kotlin, caption={KafkaConfig.kt --- default factory (file-analyzed)}]
@Bean
fun kafkaListenerContainerFactory(
    errorHandler: DefaultErrorHandler
): ConcurrentKafkaListenerContainerFactory<String, String> {
    val factory = ConcurrentKafkaListenerContainerFactory<String, String>()
    factory.consumerFactory = consumerFactory()
    factory.containerProperties.ackMode = ContainerProperties.AckMode.MANUAL
    // traceparent header present in file-analyzed -> correct child span
    factory.containerProperties.isObservationEnabled = true
    factory.setCommonErrorHandler(errorHandler)
    return factory
}
\end{lstlisting}

\begin{lstlisting}[style=kotlin, caption={KafkaConfig.kt --- Debezium factory (file-outbox)}]
@Bean
fun debeziumKafkaListenerContainerFactory(
    errorHandler: DefaultErrorHandler
): ConcurrentKafkaListenerContainerFactory<String, String> {
    val factory = ConcurrentKafkaListenerContainerFactory<String, String>()
    factory.consumerFactory = consumerFactory()
    factory.containerProperties.ackMode = ContainerProperties.AckMode.MANUAL
    // Must be explicitly false: Spring Boot auto-config sets it to true on
    // ALL factory beans if tracing is enabled.  Debezium messages carry no
    // Kafka headers, so observation would generate orphan root spans.
    factory.containerProperties.isObservationEnabled = false
    factory.setCommonErrorHandler(errorHandler)
    return factory
}
\end{lstlisting}

\texttt{MetadataProcessor} opts into the Debezium factory:
\begin{lstlisting}[style=kotlin, caption={MetadataProcessor.kt --- @KafkaListener}]
@KafkaListener(
    id = "Ai-service-consumer",
    topics = ["file-outbox"],
    containerFactory = "debeziumKafkaListenerContainerFactory"
)
override fun onMessage(message: String, acknowledgment: Acknowledgment) { ... }
\end{lstlisting}

\textbf{Important}: simply omitting \texttt{isObservationEnabled} is not
sufficient.  Spring Boot 3.x contributes a
\texttt{ConcurrentKafkaListenerContainerFactoryCustomizer} that enables
observation on every factory bean in the context when tracing is active.  The
flag must be \textbf{set to \texttt{false} explicitly} to override this.

% -----------------------------------------------------------------------------
\section{Trace Context Propagation Through the Debezium Outbox}
\label{sec:debezium}

\subsection{Architecture of the message flow}

File events do \textbf{not} go through a \texttt{KafkaTemplate}.  The actual
flow is:

\begin{enumerate}
  \item \textbf{TusUploaderService} --- \texttt{TusService}/\texttt{FileService}
        writes a row into the \texttt{files\_outbox\_events} PostgreSQL table.
  \item \textbf{Debezium} (Kafka Connect) watches the table with CDC and
        publishes each row to the \texttt{file-outbox} Kafka topic.
        Crucially, Debezium never sets Kafka record headers --- the JSON body
        is the only carrier available.
  \item \textbf{AIProcessingService} --- \texttt{MetadataProcessor} consumes
        from \texttt{file-outbox}, performs AI analysis, and writes a result
        row into \texttt{job\_outbox}.
  \item Debezium publishes \texttt{job\_outbox} rows to the
        \texttt{job-outbox} Kafka topic.
  \item \textbf{TusUploaderService} --- \texttt{KafkaConsumer} consumes from
        \texttt{job-outbox} and updates the file metadata in the database.
\end{enumerate}

\subsection{Why the RecordInterceptor approach does not work}

An earlier attempt used a \texttt{RecordInterceptor} to extract
\texttt{traceparent} from the message body and inject it into the Kafka record
headers, expecting Spring Kafka's observation to then read the header and create
a child span.

This approach \textbf{fails} because Spring Kafka creates the consumer
observation (which reads the headers) from the \emph{original}
\texttt{ConsumerRecord} \textbf{before} calling
\texttt{RecordInterceptor.intercept()}.  By the time the interceptor adds the
header, the observation has already started as a new root span.

\subsection{Correct approach: manual Propagator.extract() in the listener}

The fix is to bypass the Spring Kafka observation entirely for Debezium
consumers and manually resume the upstream trace inside the listener method,
using Micrometer's \texttt{Propagator} (auto-configured by
\texttt{micrometer-tracing-bridge-otel} as W3C TraceContext):

\begin{lstlisting}[style=kotlin, caption={Pattern used in both Debezium listeners}]
val traceparent = extractTraceparentFromPayload(message)
val span = if (traceparent != null) {
    propagator.extract(mapOf("traceparent" to traceparent)) { map, key -> map[key] }
        .name("kafka receive <topic>")
        .start()
} else {
    tracer.nextSpan().name("kafka receive <topic>").start()
}
val scope = tracer.withSpan(span)
try {
    // ... listener logic ...
} finally {
    scope.close()
    span.end()
}
\end{lstlisting}

\begin{lstlisting}[style=kotlin, caption={Debezium payload traceparent extraction}]
private fun extractTraceparentFromPayload(message: String): String? = try {
    val root = objectMapper.readTree(message)
    // Debezium serialises the JSONB payload column as an escaped JSON string.
    // Structure: { "event_type": "...", "payload": "{\"traceparent\":\"00-...\"}", ... }
    root["payload"]?.asText()?.let { payloadStr ->
        objectMapper.readTree(payloadStr)["traceparent"]
            ?.asText()?.takeIf { it.isNotBlank() }
    } ?: root["traceparent"]?.asText()?.takeIf { it.isNotBlank() }
} catch (_: Exception) { null }
\end{lstlisting}

This pattern is applied in:
\begin{itemize}[noitemsep]
  \item \textbf{AIProcessingService} --- \texttt{MetadataProcessor.onMessage()}
        (consumes \texttt{file-outbox})
  \item \textbf{TusUploaderService} --- \texttt{KafkaConsumer.consumeFileUploadCompleted()}
        (consumes \texttt{job-outbox})
\end{itemize}

\subsection{Producer side: embedding traceparent in outbox payloads}

For the context to be available to consumers, the W3C \texttt{traceparent}
string is embedded in every outbox event payload at write time.

\subsubsection{TusUploaderService --- TusService.kt and FileService.kt}

\texttt{Tracer} was injected into both services.  When building the outbox
event:

\begin{lstlisting}[style=kotlin, caption={FILE\_INSERTED / FILE\_DELETED outbox payload}]
val payload = mapOf(
    // ... existing fields ...
    "traceparent" to tracer.currentSpan()?.context()
        ?.let { "00-${it.traceId()}-${it.spanId()}-01" }
)
\end{lstlisting}

\subsubsection{AIProcessingService --- JobStatusService.kt}

\texttt{JobStatusService.createOutboxEvent()} is called inside the
\texttt{MetadataProcessor} span (which is itself a child of the upload trace
after the fix above).  \texttt{tracer.currentSpan()} therefore captures the
correct parent context for the \texttt{job-outbox} events:

\begin{lstlisting}[style=kotlin, caption={JobStatusService.kt --- createOutboxEvent()}]
private fun createOutboxEvent(event_type: String, payload: Map<String, Any>) {
    val enrichedPayload = payload.toMutableMap().apply {
        tracer.currentSpan()?.context()?.let {
            put("traceparent", "00-${it.traceId()}-${it.spanId()}-01")
        }
    }
    jobOutboxRepository.save(JobOutbox(event_type = event_type, payload = enrichedPayload))
}
\end{lstlisting}

\subsection{Producer side: KafkaTemplate messages (file-analyzed, embeddings-batch)}

Topics published via \texttt{KafkaTemplate} support explicit header injection
using \texttt{ProducerRecord}:

\begin{lstlisting}[style=kotlin, caption={MetadataProcessor.kt --- publishFileAnalyzedEvent()}]
val headers = RecordHeaders()
tracer.currentSpan()?.context()?.let {
    headers.add("traceparent",
        "00-${it.traceId()}-${it.spanId()}-01".toByteArray(Charsets.UTF_8))
}
kafkaTemplate.send(ProducerRecord(
    "file-analyzed", null, job.fileId.toString(), eventJson, headers))
\end{lstlisting}

The same pattern is used in \texttt{EmbeddingProcessor} for the
\texttt{embeddings-batch} topic.  Because these topics carry the
\texttt{traceparent} \emph{as a Kafka header}, \texttt{isObservationEnabled =
true} (the default factory) works correctly and no manual extraction is needed.

% -----------------------------------------------------------------------------
\section{TusUploaderService --- KafkaConsumerConfig.kt}

The \texttt{KafkaConsumerConfig} in TusUploaderService has a single consumer
(the \texttt{job-outbox} Debezium topic).  As with the Debezium factory in
AIProcessingService, \texttt{isObservationEnabled} must be explicitly
\texttt{false}:

\begin{lstlisting}[style=kotlin, caption={TusUploaderService --- KafkaConsumerConfig.kt}]
@Bean
fun kafkaListenerContainerFactory(): ConcurrentKafkaListenerContainerFactory<String, String> {
    val factory = ConcurrentKafkaListenerContainerFactory<String, String>()
    factory.consumerFactory = consumerFactory()
    // Must be explicitly false: Spring Boot auto-config enables observation on all
    // factories. job-outbox is a Debezium topic (no Kafka headers); trace context
    // is handled manually in the listener via Propagator.extract().
    factory.containerProperties.isObservationEnabled = false
    return factory
}
\end{lstlisting}

% -----------------------------------------------------------------------------
\section{updateFileMetadata --- Manual Span}

\texttt{FileService.updateFileMetadata()} is a plain internal method (no HTTP,
no Kafka) so it does not generate a span automatically.  To make the database
write visible in Tempo as a child of the \texttt{kafka receive job-outbox} span,
a manual span is created using the already-injected \texttt{Tracer}:

\begin{lstlisting}[style=kotlin, caption={FileService.kt --- updateFileMetadata()}]
override fun updateFileMetadata(id: UUID, version: Long?, metadata: Map<String, String>) {
    val span = tracer.nextSpan().name("updateFileMetadata").start()
    val scope = tracer.withSpan(span)
    try {
        // ... DB read and write logic ...
    } finally {
        scope.close()
        span.end()
    }
}
\end{lstlisting}

% -----------------------------------------------------------------------------
\section{Bug Fix --- AI Worker Cannot Download Files}

\subsection{Symptom}

\begin{lstlisting}[style=code, caption={Error log from AIProcessingService}]
ERROR g.a.services.MetadataProcessor
  --> Error calling file service for download <fileId>: 404 :
      {"detail":"User not found: <userId>",
       "instance":"/api/v1/files/<fileId>/download"}
\end{lstlisting}

\subsection{Root cause}

\texttt{FileService.downloadFile()} performed two operations unconditionally for
all callers regardless of role:

\begin{enumerate}
  \item \textbf{\texttt{getUserDetails(token)}}: extracts the user ID from the
        JWT, then calls \texttt{UserDetailService GET /api/v1/users/\{id\}}.
        The AI service authenticates as the \texttt{ai-client} Keycloak service
        account --- this account exists in Keycloak but has no row in
        \texttt{userdb}, so UserDetailService returned 404.

  \item \textbf{\texttt{getCurrentUserPermission()}}: annotated with
        \texttt{@PreAuthorize} that does not list \texttt{ROLE\_ai\_worker},
        so this call would have thrown a security exception regardless.
\end{enumerate}

\subsection{Fix --- FileService.kt}

\begin{lstlisting}[style=kotlin, caption={FileService.kt --- downloadFile()}]
val isAiWorker = SecurityContextHolder.getContext()
    .authentication.authorities
    .any { it.authority == "ROLE_ai_worker" }

if (!isAiWorker) {
    val token = authService.getAccessToken()
    val userDetails = getUserDetails(token)
    getCurrentUserPermission(id, userDetails)
}
\end{lstlisting}

The \texttt{@PreAuthorize} on the method already enforces that only
authenticated, authorized roles can reach this code.  For
\texttt{ROLE\_ai\_worker} callers the JWT authority check is the trust boundary;
an additional lookup in a service that does not register service accounts would
always fail.

% -----------------------------------------------------------------------------
\section{Summary of Files Modified}

\begin{center}
\begin{tabular}{lll}
\toprule
\textbf{Service} & \textbf{File} & \textbf{Change} \\
\midrule
All 4 & \texttt{build.gradle.kts}    & Removed OTel agent task; added Micrometer deps \\
All 4 & \texttt{application.yml}     & Added \texttt{management.tracing} + OTLP endpoint \\
All 4 & \texttt{logback-spring.xml}  & Added traceId/spanId to CONSOLE appender \\
\midrule
AIProcessingService & \texttt{conf/KafkaConfig.kt}
  & Two factories: default (obs=true) and Debezium (obs=false) \\
AIProcessingService & \texttt{services/MetadataProcessor.kt}
  & Propagator.extract() in listener; containerFactory override \\
AIProcessingService & \texttt{services/EmbeddingProcessor.kt}
  & traceparent header in embeddings-batch ProducerRecord \\
AIProcessingService & \texttt{services/JobStatusService.kt}
  & traceparent embedded in job\_outbox event payload \\
\midrule
TusUploaderService & \texttt{conf/KafkaConsumerConfig.kt}
  & isObservationEnabled = false (explicit) \\
TusUploaderService & \texttt{services/KafkaConsumer.kt}
  & Propagator + Tracer injected; Propagator.extract() in listener \\
TusUploaderService & \texttt{services/TusService.kt}
  & Tracer injected; traceparent in FILE\_INSERTED payload \\
TusUploaderService & \texttt{services/FileService.kt}
  & traceparent in FILE\_DELETED payload \\
TusUploaderService & \texttt{services/FileService.kt}
  & downloadFile: skip user lookup for ROLE\_ai\_worker \\
TusUploaderService & \texttt{services/FileService.kt}
  & updateFileMetadata: manual span added \\
\bottomrule
\end{tabular}
\end{center}

\end{document}